\begin{table}[H]
\centering
\caption{Robustness Checks}
\label{tab:robust}
\begin{threeparttable}
\begin{tabular}{lccc}
\toprule
Specification & Coefficient & SE & N \\
\midrule
\multicolumn{4}{l}{\textit{Panel A: Simple DiD (High-Heat $\times$ Post)}} \\
Baseline (2016--2023) & -0.348*** & (0.095) & 2,337,654 \\
Restricted: 2019--2023 & 0.021 & (0.100) & 1,594,391 \\
Excluding 2020--2021 & -0.744*** & (0.104) & 1,744,971 \\
Federal OSHA states only & -0.271*** & (0.069) & 1,361,151 \\
State-plan states (placebo) & -0.463** & (0.181) & 976,503 \\
Large establishments ($\geq$250 empl.) & -0.248 & (0.164) & 234,304 \\
\addlinespace
\multicolumn{4}{l}{\textit{Panel B: Triple DiD (High-Heat $\times$ Post $\times$ Hot State)}} \\
Baseline (2016--2023) & 0.046 & (0.168) & 2,337,653 \\
Restricted: 2019--2023 & 0.003 & (0.177) & 1,594,388 \\
Federal OSHA states only & 0.259** & (0.105) & 1,361,150 \\
Large establishments ($\geq$250 empl.) & -0.067 & (0.268) & 234,279 \\
\bottomrule
\end{tabular}
\begin{tablenotes}[flushleft]
\small
\item \textit{Notes:} Panel A reports the simple DiD coefficient (High-Heat $\times$ Post) under alternative samples. Panel B reports the triple-DiD coefficient. All specifications include state$\times$year fixed effects. Panel B specifications add industry$\times$state fixed effects. The ``Restricted: 2019--2023'' specification uses a shorter pre-treatment window that produces near-zero estimates, confirming that the full-sample DiD result reflects secular trends rather than the NEP. The state-plan placebo shows a significant ``effect'' in jurisdictions where the federal NEP does not apply. Standard errors clustered at the state level. * $p<0.10$, ** $p<0.05$, *** $p<0.01$.
\end{tablenotes}
\end{threeparttable}
\end{table}

